\documentclass[11pt,a4paper]{article}
\usepackage[margin=2.0cm]{geometry}
\usepackage[T1]{fontenc}
\usepackage{lmodern}
\usepackage{amsmath}
\usepackage{booktabs}
\usepackage{tabularx}
\usepackage{xcolor}
\usepackage{listings}
\usepackage[hidelinks]{hyperref}

\emergencystretch=2em
\definecolor{codebg}{gray}{0.95}
\lstset{language=C++, basicstyle=\ttfamily\footnotesize, keywordstyle=\color{blue!70!black},
  commentstyle=\color{green!50!black}\itshape, backgroundcolor=\color{codebg},
  frame=single, rulecolor=\color{black!20}, breaklines=true, showstringspaces=false,
  tabsize=2, columns=flexible}
\newcommand{\code}[1]{\texttt{#1}}

\title{ATmega32A Module \\ \Large Timers (Normal, CTC and Fast PWM)}
\author{Ali Sahafi}
\date{}

\begin{document}
\maketitle

\section{Theory}

A hardware timer is a counter. It counts clock ticks automatically, and your
program does not need to do anything while it counts. The ATmega32A has three
timers:

\begin{center}
\begin{tabular}{lllll}
\toprule
\textbf{Timer} & \textbf{Size} & \textbf{Counts} & \textbf{Prescalers} & \textbf{OC pin(s)} \\
\midrule
\code{Timer0} & 8-bit  & 0--255   & 1, 8, 64, 256, 1024          & \code{PB3} (LED D3) \\
\code{Timer1} & 16-bit & 0--65535 & 1, 8, 64, 256, 1024          & \code{PD5} (A), \code{PD4} (B) \\
\code{Timer2} & 8-bit  & 0--255   & 1, 8, 32, 64, 128, 256, 1024 & \code{PD7} \\
\bottomrule
\end{tabular}
\end{center}

\paragraph{Prescaler.} The prescaler $N$ divides the CPU clock before it goes
to the counter. So one count takes $t_{\text{count}} = N / F_{\text{CPU}}$.
Example: with 8\,MHz and $N = 1024$, one count takes $128\,\mu$s.

\paragraph{The three modes.}
\begin{itemize}
  \item \textbf{Normal} (\code{TIMER\_NORMAL}): The timer counts up to the
        maximum (255, or 65535 for Timer1). Then it goes back to 0. This is
        called an \emph{overflow}, and the timer sets the overflow flag
        \code{TOV}. For a shorter time, you give the counter a start value
        (a \emph{preload}):
        \[ t = (256 - \text{preload}) \cdot \frac{N}{F_{\text{CPU}}} \qquad (\text{use } 65536 \text{ for Timer1}) \]
  \item \textbf{CTC} (\code{TIMER\_CTC}, \emph{Clear Timer on Compare}):
        The timer counts up to the compare value \code{OCR}. Then it sets the
        compare flag \code{OCF} and starts again from 0. You do not need a
        preload:
        \[ t = (\text{OCR} + 1) \cdot \frac{N}{F_{\text{CPU}}} \]
        On every compare match, the timer can also toggle, set or clear its
        OC pin. This gives a square wave, and the CPU does not need to do anything.
  \item \textbf{Fast PWM} (\code{TIMER\_FAST\_PWM}): The timer counts from 0 to
        255 again and again. The OC pin changes at 0 and at the compare value.
        The result is a PWM signal. \code{OCR} sets the duty cycle:
        \[ F_{\text{PWM}} = \frac{F_{\text{CPU}}}{N \cdot 256} \qquad
           \text{Duty cycle (active part)} = \frac{\text{OCR} + 1}{256} \]
\end{itemize}

\section{How to use the driver}

The driver gives you three objects: \code{Timer0}, \code{Timer1} and
\code{Timer2}. You start a timer with one function call:

\begin{lstlisting}
Timer0.begin(mode, prescaler);            // start the timer
Timer0.begin(mode, prescaler, pinMode);   // start the timer and use the OC pin
\end{lstlisting}

\subsection*{The options}

\begin{center}
\small
\begin{tabularx}{\textwidth}{@{}l l X@{}}
\toprule
\textbf{Argument} & \textbf{Value} & \textbf{Meaning} \\
\midrule
\code{mode} & \code{TIMER\_NORMAL}   & Count from 0 to the maximum, overflow, start again from 0 \\
            & \code{TIMER\_CTC}      & Count from 0 to the compare value, start again from 0 \\
            & \code{TIMER\_FAST\_PWM} & Count from 0 to 255, the OC pin gives a PWM signal \\
\midrule
\code{prescaler} & \code{1}, \code{8}, \code{64}, \code{256}, \code{1024} &
      Write the number you want to divide by. Timer2 can also use \code{32} and \code{128}. \\
      & \code{EXT\_FALLING}, \code{EXT\_RISING} &
      Only Timer0 and Timer1: count pulses on pin \code{T0}\,=\,\code{PB0} or \code{T1}\,=\,\code{PB1} (counter mode) \\
\midrule
\code{pinMode} & \code{OC\_OFF} (default) & The OC pin is a normal GPIO pin \\
\multicolumn{3}{@{}l}{\emph{\quad with \code{TIMER\_NORMAL} or \code{TIMER\_CTC}:}} \\
               & \code{OC\_TOGGLE} & Toggle the OC pin on every compare match \\
               & \code{OC\_CLEAR}  & Set the OC pin LOW on a compare match \\
               & \code{OC\_SET}    & Set the OC pin HIGH on a compare match \\
\multicolumn{3}{@{}l}{\emph{\quad with \code{TIMER\_FAST\_PWM}:}} \\
               & \code{PWM\_NON\_INVERTING} & HIGH from 0 to the compare value, then LOW \\
               & \code{PWM\_INVERTING}      & LOW from 0 to the compare value, then HIGH.
                  Use this for the \textbf{active-low LEDs} on the board. \\
\bottomrule
\end{tabularx}
\end{center}

If \code{pinMode} is not \code{OC\_OFF}, the driver makes the OC pin an output
for you. While the timer controls this pin, \code{GPIO.write()} has no effect
on it.

\subsection*{Functions}

\begin{center}
\small
\begin{tabularx}{\textwidth}{@{}l X@{}}
\toprule
\textbf{Function} & \textbf{What it does} \\
\midrule
\code{Timer0.begin(mode, prescaler, pinMode)} & Set up and start the timer (see the options above) \\
\code{Timer0.setCount(value)} & Write a start value to the counter (example: 156 gives 100 counts in Normal mode) \\
\code{Timer0.getCount()} & Read the counter (0--255; Timer1: 0--65535) \\
\code{Timer0.setCompare(value)} & Set the compare value. CTC: the timer counts up to this value. PWM: the level, 0--255 \\
\code{Timer0.setDutyCycle(percent)} & Only PWM: how many percent (0--100) of the time the pin is \emph{active}.
       Active is HIGH for \code{PWM\_NON\_INVERTING} and LOW for \code{PWM\_INVERTING}. 0\,\% and 100\,\% are exact. \\
\code{Timer0.overflowed()} & Returns \code{true} one time for each overflow. The flag is cleared for you \\
\code{Timer0.compareMatched()} & Returns \code{true} one time for each compare match. The flag is cleared for you \\
\code{Timer0.stop()} / \code{Timer0.start()} & Stop / continue counting. The settings do not change \\
\code{Timer0.onOverflow(function)} & \emph{Optional:} call \code{function} on every overflow (interrupt) \\
\code{Timer0.onCompareMatch(function)} & \emph{Optional:} call \code{function} on every compare match (interrupt) \\
\midrule
\multicolumn{2}{@{}l}{\emph{Only Timer1 (it has two OC pins: A = \code{PD5}, B = \code{PD4}):}} \\
\code{Timer1.begin(mode, N, pinModeA, pinModeB)} & Channel B has its own \code{pinMode} \\
\code{Timer1.setCompareA/B(value)} & 16-bit compare values. \code{setCompare()} is channel A. In CTC mode the timer counts up to A \\
\code{Timer1.setDutyCycleA/B(percent)} & PWM duty cycle for each channel. \code{setDutyCycle()} is channel A \\
\code{Timer1.compareMatchedA/B()} & Compare flag for each channel. \code{compareMatched()} is channel A \\
\code{Timer1.setTop(top)} & Length of the PWM period in counts (default 255): $T_{\text{PWM}} = (\text{top}+1) \cdot N / F_{\text{CPU}}$ \\
\bottomrule
\end{tabularx}
\end{center}

\paragraph{Waiting for a flag (polling).}
\code{overflowed()} and \code{compareMatched()} return \code{true} when the event
happened after your last call. So \code{while (!Timer0.overflowed());} waits
until the next overflow. You do not need to write 1 to \code{TIFR}.

\paragraph{Interrupts (optional).} \code{onOverflow()} and
\code{onCompareMatch()} turn on the interrupt and call \code{sei()} for you.
Your function then runs automatically. Keep this function short. If it
changes a variable that \code{main()} also uses, declare the variable as
\code{volatile}. To turn the interrupt off again, pass \code{nullptr}.

\subsection*{Driver functions and registers}

In the lecture slides we set the registers by hand. Each driver function does
one of these steps for you:

\begin{center}
\small
{\footnotesize\begin{tabular}{@{}lll@{}}
\toprule
\textbf{Driver function} & \textbf{Registers (Timer0)} & \textbf{Bits} \\
\midrule
\code{Timer0.begin(TIMER\_NORMAL, 1024)}          & \code{TCCR0 = 0x05;} & \code{CS02 CS00} \\
\code{Timer0.begin(TIMER\_CTC, 1024, OC\_TOGGLE)} & \code{TCCR0 = 0x1D;} & \code{WGM01 COM00 CS02 CS00} \\
                                                  & \code{DDRB |= 1<{}<PB3;} & OC0 pin as output \\
\code{Timer0.begin(TIMER\_FAST\_PWM, 1, PWM\_INVERTING)} & \code{TCCR0 = 0x79;} & \code{WGM00 WGM01 COM01 COM00 CS00} \\
\code{Timer0.setCount(156)}   & \code{TCNT0 = 156;} & \\
\code{Timer0.setCompare(x)}   & \code{OCR0 = x;} & \\
\code{Timer0.overflowed()}    & \code{TIFR \& (1<{}<TOV0)} & then \code{TIFR = 1<{}<TOV0;} \\
\code{Timer0.compareMatched()} & \code{TIFR \& (1<{}<OCF0)} & then \code{TIFR = 1<{}<OCF0;} \\
\bottomrule
\end{tabular}}
\end{center}

\paragraph{Important: \code{F\_CPU} and the fuses.} All timing formulas use the
\emph{real} clock of the chip. \code{make} and \code{make timer} always set the
fuses for the external 8\,MHz crystal. To use another clock:
(1) set the fuses, for example \code{make int-1mhz};
(2) upload your program without changing the fuses:
\code{make build flash SRC=Task.cpp};
(3) change \code{\#define F\_CPU} to the new clock, so that
\code{\_delay\_ms()} is still correct.

\newpage
\section{Example: blinking LED and PWM, both done by the timers}

\paragraph{Board hardware.} All parts are already on the development board:
\begin{itemize}
  \item \textbf{8 LEDs (D0--D7) on PORTB}. They are active-low: \code{LOW} = ON. LED \textbf{D3} is on \code{PB3} = \code{OC0}, the output pin of Timer0.
  \item \textbf{Seven-segment display}, also on PORTB (it uses the same pins as the LEDs).
  \item \textbf{8 DIP switches (S0--S7) on PORTC}. A closed switch reads \code{LOW}. Use \code{INPUT\_PULLUP}.
  \item \textbf{Push buttons}: S11 on \code{PD2} and S10 on \code{PD6}. A pressed button reads \code{LOW}.
\end{itemize}

\paragraph{Build and upload.} \code{make timer}

\lstinputlisting{example.cpp}

\newpage
\section{Laboratory Tasks}

For each task, make a new \code{.cpp} file in the main project folder (for
example \code{Task1.cpp}). Start the file with \code{\#define F\_CPU}. Then
include the drivers you need: \code{drivers/gpio/gpio.hpp},
\code{drivers/timer/timer.hpp}, \code{drivers/common/sevenseg.hpp}. Upload it
with \code{make SRC=Task1.cpp}.

\subsection*{Task 0: Fuse Bits and \code{F\_CPU} (10 min)}
Show the numbers 0 to 9 on the seven-segment display. Change the number every
second. Use \code{SevenSeg(digit)} and \code{\_delay\_ms(1000)}.
\begin{itemize}
  \item \textbf{Task 0.1 (Internal 1\,MHz):} Run \code{make int-1mhz}. Write
        \code{\#define F\_CPU 1000000UL} and upload with
        \code{make build flash SRC=Task0.cpp}. Check the time with a stopwatch.
  \item \textbf{Task 0.2 (External 8\,MHz crystal):} Write
        \code{\#define F\_CPU 8000000UL} and upload with \code{make SRC=Task0.cpp}.
        This command also selects the crystal.
  \item \textbf{Task 0.3 (Wrong \code{F\_CPU}):} Keep \code{F\_CPU 8000000UL}, but
        change the chip back to 1\,MHz with \code{make int-1mhz}. How long is one
        ``second'' now? Explain why.
\end{itemize}

\subsection*{Task 1: Timer Normal Mode --- Counter on the LEDs (30 min)}
Count in binary on the 8 LEDs. Add 1 every \textbf{100\,ms}. Use \code{Timer0}
in \code{TIMER\_NORMAL} mode for the timing (external 8\,MHz crystal). The LEDs
are active-low, so write the inverted value:\\ \code{GPIO.write(PORTB, \textasciitilde count)}.
\begin{itemize}
  \item \textbf{Task 1.1 (Calculate first):} With prescaler 1024, how long is one
        count? What is the \emph{longest} time between two Timer0 overflows?
        Why is one overflow not enough for 100\,ms?
  \item \textbf{Task 1.2 (10 $\times$ 10\,ms):} Calculate the preload value for
        10\,ms. In the main loop, do these two steps 10 times:
        \begin{enumerate}
          \item write the preload with \code{Timer0.setCount()};
          \item wait with \code{while (!Timer0.overflowed());}.
        \end{enumerate}
        After that, add 1 to the counter.
  \item \textbf{Task 1.3 (16-bit Timer1):} Do the same with \code{Timer1}.
        Timer1 can make 100\,ms with \emph{one} overflow. With which prescalers
        is the preload value a whole number? Compare your code with Task 1.2.
\end{itemize}

\subsection*{Task 2: Timer CTC Mode --- Blinking in Hardware (20 min)}
The timer blinks LED \textbf{D3} (\code{PB3 = OC0}) without help from the CPU.
The DIP switches on \code{PORTC} set the blink speed.
\begin{itemize}
  \item \textbf{Task 2.1 (Blink with the timer):} Start \code{Timer0} in
        \code{TIMER\_CTC} mode with prescaler 1024 and \code{OC\_TOGGLE}. The main
        loop only does one thing: it reads the switches (\code{GPIO.read(PINC)})
        and writes the value with \code{Timer0.setCompare()}. You do not need to
        check any flag. At 8\,MHz, calculate the slowest possible blink frequency.
        What do you see on D3? Explain why.
  \item \textbf{Task 2.2 (Make it visible):} Change the chip to the internal
        1\,MHz clock (see Task 0.1). What is the shortest and the longest time
        between two toggles now? Test it with the switches.
  \item \textbf{Task 2.3 (Two things at the same time):} Keep the blinking LED.
        In the same loop, also show the switch value on the other LEDs:
        \code{GPIO.write(PORTB, GPIO.read(PINC))}. The loop also writes to
        \code{PB3}. Why does D3 still blink?
\end{itemize}

\subsection*{Task 3: Timer Fast PWM --- LED Brightness (25 min)}
Change the brightness of LED \textbf{D3} (\code{OC0}) with PWM from \code{Timer0}.
\begin{itemize}
  \item \textbf{Task 3.1 (Switches $\rightarrow$ brightness):} Start \code{Timer0}
        in \code{TIMER\_FAST\_PWM} mode with \code{PWM\_INVERTING}, because the
        LED is active-low. Choose the prescaler so that
        $F_{\text{PWM}} \approx 4$\,kHz at 8\,MHz. In the loop, write
        \code{Timer0.setCompare(GPIO.read(PINC))}.
  \item \textbf{Task 3.2 (Buttons, in percent):} Use the push buttons instead of
        the switches. S11 (\code{PD2}) makes the LED 10\,\% brighter. S10
        (\code{PD6}) makes it 10\,\% darker. Keep the value between 0 and
        100\,\%. Use \code{Timer0.setDutyCycle(percent)}. One press must change
        the value only one time, so wait until the button is released.
  \item \textbf{Task 3.3 (Flicker):} Calculate $F_{\text{PWM}}$ for the prescalers
        64, 256 and 1024. Try all three. With which prescaler do you start to
        see the LED flicker?
  \item \textbf{Task 3.4 (Breathing LED with two timers):} Make D3 slowly brighter
        from 0 to 100\,\%, then slowly darker again. Change the value by 1\,\%
        every 10\,ms. Do not use \code{\_delay\_ms()}. Instead, use \code{Timer1}
        in \code{TIMER\_CTC} mode to measure the 10\,ms
        (\code{Timer1.compareMatched()}). Calculate the prescaler and the compare
        value for Timer1.
\end{itemize}

\end{document}
